\documentclass[11pt]{article}
\usepackage{amsmath,amssymb,amsthm}
\newtheorem{proposition}{Proposition}
\title{A Variable-Pole Transition in Compact Reciprocal Permutation Sheets}
\author{}
\date{November 23, 2026}
\begin{document}
\maketitle

\paragraph{Scope.}
This note checks one structural possibility in the joint reverse compiler.
It proves that varying the common row pole can change the affine order
type of permutation-sheet values, even arbitrarily far from all token
poles. Thus variable-pole queries cannot simply be identified with one
fixed linear embedding of the permutations. It proves no shattering
lower bound and no superconstant head saving. All statements below have
written proofs and an exact symbolic checker, but are not Lean formalized.

For four distinct real tokens $\lambda_1,\ldots,\lambda_4$ and
$d>\max_j\lambda_j$, associate to a permutation $\pi$ the point
\[
 v_\pi(d)=\left(\frac1{d-\lambda_{\pi(1)}},
                   \frac1{d-\lambda_{\pi(2)}},
                   \frac1{d-\lambda_{\pi(3)}}\right)\in\mathbb R^3.
\]
Omitting the fourth coordinate loses no affine information because the
sum of all four coordinates is independent of $\pi$.
Use the four permutations, in one-based notation,
\[
 \pi_1=(1,2,3,4),\quad \pi_2=(1,3,4,2),\quad
 \pi_3=(2,4,3,1),\quad \pi_4=(4,1,2,3).
\]
Let $\Delta(d)$ be the determinant whose three rows are
$v_{\pi_i}(d)-v_{\pi_4}(d)$, $1\leq i\leq3$.

\begin{proposition}
The sign of $\Delta(d)$ need not be constant for
$d>\max_j\lambda_j$. Moreover, positive tokens in $[1,2]$ can be
chosen so that a sign-changing zero occurs beyond any prescribed bound.
\end{proposition}
\begin{proof}
Put $P(d)=\prod_{j=1}^4(d-\lambda_j)>0$.
Replacing each reciprocal by
$P(d)/(d-\lambda_j)=\prod_{a\neq j}(d-\lambda_a)$ shows that
$Q(d)=P(d)^3\Delta(d)$ is a polynomial. For
$(\lambda_1,\lambda_2,\lambda_3,\lambda_4)=(112,273,707,788)$,
exact expansion gives
\[
\begin{split}
Q(d)={}&-5(d-707)(d-112)\bigl(5275d^4-6245360632d^3\\
 &\qquad+8524750604517d^2-3261701068819508d
                       +248767066754926160\bigr).
\end{split}
\]
Its values at $10^4$ and $10^8$ are respectively positive and negative.
Both arguments exceed every token, so the same is true of $\Delta$.

For arbitrary distance use tokens $(0,1/4,c,1)$, where $1/4<c<1$.
The polynomial $Q_c$ has degree at most six, with leading coefficients
\[
 [d^6]Q_c=-\frac{16c^2+44c-51}{64}=L(c),\qquad
 [d^5]Q_c=\frac{c^3}{4}+\frac{9c^2}{4}
                       -\frac{47c}{32}-\frac{15}{32}=N(c).
\]
Let $c_0=(-11+5\sqrt{13})/8\in(1/4,1)$.
Then $L(c_0)=0$ and
\[
 N(c_0)=\frac{363}{32}-\frac{795\sqrt{13}}{256}>0,
\]
where positivity follows from $2904^2>795^2\cdot13$.
Consequently $Q_{c_0}(d)>0$ for every sufficiently large $d$.
Given a prescribed bound $M$, choose $D>M$ sufficiently large that
$Q_{c_0}(D)>0$. Continuity in $c$ permits choosing $c>c_0$ so close
that $Q_c(D)>0$. Since $L'(c)=-c/2-11/16<0$, this choice has
$L(c)<0$, and hence $Q_c(d)<0$ for all sufficiently large $d$.
There is therefore a sign-changing root beyond $D$.
Adding one to every token and to the pole places the tokens in
$[1,2]$ without changing any reciprocal difference or determinant.
\end{proof}

\paragraph{Compatibility with full-cube positivity.}
This transition is not forced to occur near a forbidden pole.
For any fixed number of rows $r$ and columns $s$, take a sign-changing
root $d_0>4s+4$, after the preceding translation. There are opposite
signs at poles arbitrarily close to $d_0$: a polynomial has finitely
many roots, and a sign-changing root has odd multiplicity. Choose
all desired row poles $d_i$ in a sufficiently short interval around
$d_0$, and choose a common intercept $C>\max_i d_i$ in that interval.
Then every row coefficient $p_i=C-d_i$ is strictly positive and
their sum can be made less than $1$. For each head and column let
the column coefficient be one of the four tokens, so it lies in
$[1,2]$. The affine denominator
\[
 B_h(z,w)=C-\sum_i p_i z_i-\sum_j u_{hj}w_j
\]
has strictly negative slopes and cube minimum at least
$C-1-2s>0$. Thus the orientation transition itself is compatible
with the full-cube atom restriction. This observation does not show
that any additional family of signs is independently decodable.

\paragraph{Remaining gap.}
For one fixed pole all permutation points lie in one three-dimensional
affine space, and four affine tests cannot shatter sixteen patterns.
The transition proves only that the family of available affine cuts
changes when the pole changes. Combining such changes into even a
four-bit decoder, and more importantly into an unbounded row/head
ratio, requires a separate argument. Disjoint repetition of finite
transitions gives at most a constant ratio and does not solve the
requested asymptotic problem.
\end{document}
